
\documentclass[
	% -- opções da classe memoir --
	12pt,				% tamanho da fonte
	% openright,			% capítulos começam em pág ímpar (insere página vazia caso preciso)
	oneside,			% para impressão apenas no anverso (apenas frente). Oposto a twoside
	a4paper,			% tamanho do papel. 
	% -- opções da classe abntex2 --
	%chapter=TITLE,		% títulos de capítulos convertidos em letras maiúsculas
	%section=TITLE,		% títulos de seções convertidos em letras maiúsculas
	%subsection=TITLE,	% títulos de subseções convertidos em letras maiúsculas
	%subsubsection=TITLE,% títulos de subsubseções convertidos em letras maiúsculas
	% -- opções do pacote babel --
	english,			% idioma adicional para hifenização
	%french,				% idioma adicional para hifenização
	%spanish,			% idioma adicional para hifenização
	brazil				% o último idioma é o principal do documento
	]{abntex2ppgsi}

% ---
% Pacotes básicos 
% ---
% \usepackage{lmodern}			% Usa a fonte Latin Modern			
% \usepackage[T1]{fontenc}		% Selecao de codigos de fonte.
\usepackage[utf8]{inputenc}		% Codificacao do documento (conversão automática dos acentos)
\usepackage{lastpage}			% Usado pela Ficha catalográfica
\usepackage{indentfirst}		% Indenta o primeiro parágrafo de cada seção.
\usepackage{color}				% Controle das cores
\usepackage{graphicx}			% Inclusão de gráficos
\usepackage{microtype} 			% para melhorias de justificação
\usepackage{pdfpages}     %para incluir pdf
\usepackage{algorithm}			%para ilustrações do tipo algoritmo
\usepackage{mdwlist}			%para itens com espaço padrão da abnt
\usepackage[noend]{algpseudocode}			%para ilustrações do tipo algoritmo
		
% ---
% Pacotes adicionais, usados apenas no âmbito do Modelo Canônico do abnteX2
% ---
\usepackage{lipsum}				% para geração de dummy text
% ---

% ---
% Pacotes de citações
% ---
\usepackage{hyperref}
\usepackage[brazilian,hyperpageref]{backref}	 % Paginas com as citações na bibl
\usepackage[alf,abnt-etal-list=0,abnt-etal-text=it]{abntex2cite}	% Citações padrão ABNT
\usepackage{multirow}
% para tabelas melhores 
\usepackage{tikz}
% para figuras melhores 

\usetikzlibrary{positioning}

\renewcommand{\backrefpagesname}{Citado na(s) página(s):~}
% Texto padrão antes do número das páginas
\renewcommand{\backref}{}
% Define os textos da citação
\renewcommand*{\backrefalt}[4]{
	\ifcase #1 %
		Nenhuma citação no texto.%
	\or
		Citado na página #2.%
	\else
		Citado #1 vezes nas páginas #2.%
	\fi}%

\instituicao{
	UNIVERSIDADE DE SÃO PAULO
	\par
	ESCOLA DE ARTES, CIÊNCIAS E HUMANIDADES
	\par
	PROGRAMA DE PÓS-GRADUAÇÃO EM SISTEMAS DE INFORMAÇÃO}

\titulo{Agentes Inteligentes baseados em Large Language Models para apoio à tomada de decisão em sistemas de monitoramento da qualidade da água}

\autor{\uppercase{Arthur Garcia Takahashi}}

\local{São Paulo}

\data{2027}
\orientador{Prof. Dr. Marcelo Lauretto}

\coorientador{Prof. Dr. Fábio Nakano}

\tipotrabalho{Dissertação (Mestrado) }

\definecolor{blue}{RGB}{41,5,195}

% informações do PDF
\makeatletter
\pdfstringdefDisableCommands{%
  \let\uppercase\relax
  \let\MakeUppercase\relax
  \let\newline\space
  \let\\\space
}
\hypersetup{
     	%pagebackref=true,
		pdftitle={\@title}, 
		pdfauthor={\@author},
    	pdfsubject={\imprimirpreambulo},
	    pdfcreator={laTeX com abnTeX2 adaptado para o PPgSI-EACH-USP},
		pdfkeywords={abnt}{latex}{abntex}{abntex2ppgsi}{qualificação de mestrado}{dissertação de mestrado}{qualificação de doutorado}{tese de doutorado}{ppgsi}, 
		colorlinks=true,       		% false: boxed links; true: colored links
    	linkcolor=blue,          	% color of internal links
    	citecolor=blue,        		% color of links to bibliography
    	filecolor=magenta,      		% color of file links
		urlcolor=blue,
		bookmarksdepth=4
}
\makeatother
% --- 

% --- 
% Espaçamentos entre linhas e parágrafos 
% --- 

% O tamanho do parágrafo é dado por:
\setlength{\parindent}{1.25cm}

% Controle do espaçamento entre um parágrafo e outro:
\setlength{\parskip}{0cm}  % tente também \onelineskip
\renewcommand{\baselinestretch}{1.5}

% ---
% compila o indice
% ---
\makeindex
% ---

	% Controlar linhas orfas e viuvas
  \clubpenalty10000
  \widowpenalty10000
  \displaywidowpenalty10000

% ----
% Início do documento
% ----

\begin{document}

% Retira espaço extra obsoleto entre as frases.
\frenchspacing 

% ----------------------------------------------------------
% ELEMENTOS PRÉ-TEXTUAIS
% ----------------------------------------------------------
% \pretextual

% ---
% Capa
% ---
%-------------------------------------------------------------------------
% Comentário adicional do PPgSI - Informações sobre a ``capa'':
%
% Esta é a ``capa'' principal/oficial do trabalho, a ser impressa apenas 
% para os casos de encadernação simples (ou seja, em ``espiral'' com 
% plástico na frente).
% 
% Não imprimir esta ``capa'' quando houver ``capa dura'' ou ``capa brochura'' 
% em que estas mesmas informações já estão presentes nela.
%
%-------------------------------------------------------------------------
\imprimircapa
% ---

% ---
% Folha de rosto
% (o * indica que haverá a ficha bibliográfica)
% ---
\imprimirfolhaderosto*

% ---

% ---
% RESUMOS
% ---

% resumo em português
\setlength{\absparsep}{18pt} % ajusta o espaçamento dos parágrafos do resumo


% ---
% inserir o sumario
% ---
\pdfbookmark[0]{\contentsname}{toc}
\tableofcontents*
\cleardoublepage
% ---




% ----------------------------------------------------------
% ELEMENTOS TEXTUAIS
% ----------------------------------------------------------
\textual
\chapter{Introdução}
\chapter{Fundamentação teórica}
\label{cap:fundamentacao}

\section{Inteligência Artificial}
A inteligência Artificial não pode simplesmente ser definida por uma única frente, ela é dividida por \citeonline{russell2021artificial} em quatro abordagens baseadas em duas dimensões conceituais que são Humano versus Racional e Pensamento versus Comportamento.
A partir dessas dimensões, devivam-se as quatro aborgadens: Agir como humanos (Teste de Turing), Pensar como humanos (Modelagem Cognitiva),
Pensar racionalmente (Leis de Lógica) e Agir racionalmente (Agentes Racionais). Esta último veio a se tornar o modelo padrão e mais bem-sucedido do que podemos chamar de IA moderna por focar em agir de forma racional com base em conhecimento \cite[p. 22]{russell2021artificial}.

Uma evolução importante do aprendizado de máquina ocorreu com o
desenvolvimento das técnicas de aprendizado profundo (\textit{deep learning}).
Essas técnicas utilizam modelos compostos por múltiplas camadas de
processamento capazes de aprender representações dos dados em diferentes níveis de abstração. Dessa forma, características relevantes para a realização de uma determinada tarefa podem ser aprendidas diretamente a partir dos dados, reduzindo a necessidade de definir manualmente todas as características que serão utilizadas pelo modelo \cite{lecun2015deep}.

O aprendizado profundo apresentou avanços significativos em diferentes
domínios, incluindo reconhecimento de fala, reconhecimento e detecção de
objetos, processamento de imagens, vídeo e áudio, além de aplicações em
descoberta de fármacos e genômica. Estes avanços estão associados, entre outros
fatores, à capacidade das redes neurais profundas de aprender representações
hierárquicas a partir de grandes conjuntos de dados
\cite{lecun2015deep}.

\section{Agentes inteligentes}
\label{sec:agentes-inteligentes}

Para \citeonline{russell2021artificial} agente é qualquer entidade que percebe seu ambiente por meio de sensores e age sobre ele por meio de atuadores.

O conceito de agente inteligente constitui uma das abordagens utilizadas na Inteligência Artificial para representar sistemas capazes de interagir com um ambiente e realizar ações orientadas a objetivos previamente estabelecidos. Nesse contexto, um agente pode ser compreendido como uma entidade que recebe informações sobre o ambiente, processa essas
informações e seleciona ações a serem tomadas de acordo com os objetivos \cite{muller1999}.

Müller \cite{muller1999} organiza as principais abordagens para arquiteturas de
agentes em categorias que incluem arquiteturas reativas, deliberativas e
voltadas à interação, além de discutir arquiteturas híbridas que combinam
características dessas diferentes abordagens. Essa perspectiva permite
compreender os agentes não apenas a partir de suas capacidades individuais,
mas também considerando a forma como percebem o ambiente, tomam decisões e
interagem com outras entidades.

Além das características arquiteturais, os agentes podem ser classificados
considerando sua relação com o ambiente e seu modo de interação. Müller
\cite{muller1999} apresenta uma classificação que considera, entre outros
aspectos, o estado material do agente e seu modo primário de interação.
Quanto ao estado material, distinguem-se agentes de hardware, que possuem
presença física e interagem com ambientes físicos por meio de sensores e
atuadores, e agentes de software, que interagem com ambientes computacionais
reais ou virtuais. Quanto ao modo de interação, são considerados agentes
autônomos, sistemas multiagentes e agentes assistentes
\cite{muller1999}.

A fim de ilustrar a relação prática dos fundamentos teóricos de agentes descritos por \citeonline{muller1999}, apresenta-se um mapeamento conceitual em forma de tabela resumida. Como pode ser observado no Quadro~\ref{qua:muller_agentes}, os diferentes paradigmas arquiteturais e taxonômicos são relacionados inclusive para ilustrar de forma mais didádica foi adicionado na última coluna, a título de exemplo, um cenário hipotético de monitoramento e supervisão hídrica de alta frequência.

\begin{quadro}[htb]
\centering
\caption{Paradigmas de agentes segundo Müller (1999) e sua aplicação no sistema proposto}
\label{qua:muller_agentes}
\small
\begin{tabular}{|p{2.8cm}|p{2.8cm}|p{5.2cm}|p{4.2cm}|}
\hline
\textbf{Categoria} & 
\textbf{Classificação} & 
\textbf{Descrição Teórica} & 
\textbf{Aplicação no \newline Monitoramento Hídrico} \\ \hline

Arquitetura de \newline Controle & 
Reativa & 
Decisões baseadas em estímulos diretos, 
sem modelo interno do mundo. & 
Leitura em tempo real e processamento 
imediato de alertas básicos. \\ \cline{2-4}

& 
Deliberativa & 
Raciocínio lógico, planejamento e representações 
internas complexas. & 
Agente de IA interpretando as anomalias 
e formulando diagnósticos. \\ \cline{2-4}

& 
Híbrida & 
Combinação de camadas reativas (tempo real) 
e deliberativas (metas). & 
Arquitetura central proposta para o monitoramento 
de alta frequência. \\ \cline{2-4}

& 
De Interação & 
Foco na coordenação, cooperação e comunicação 
entre agentes. & 
Orquestração do Sistema Multiagente (SMA) 
para monitoramento hídrico. \\ \hline

Taxonomia: \newline Estado Material & 
Hardware & 
Agentes com presença física, interagindo 
via sensores físicos. & 
Sondas multiparamétricas instaladas 
\textit{in situ} nos rios. \\ \cline{2-4}

& 
Software & 
Agentes virtuais operando em ambientes 
computacionais ou digitais. & 
Agentes cognitivos baseados em LLMs 
e modelos de IA. \\ \hline

Taxonomia: \newline Modo de Interação & 
Autônomo & 
Foco na relação direta e independente entre o 
próprio agente e seu ambiente. & 
Sonda hídrica coletando e transmitindo dados 
de forma isolada por telemetria. \\ \cline{2-4}

& 
Multiagente & 
Presença de outros agentes no ambiente, com 
cooperação e compartilhamento de informações. & 
Nós e sondas cooperando para analisar a 
qualidade da água ao longo de toda a bacia. \\ \cline{2-4}

& 
Assistente & 
Projetados para interagir diretamente com 
humanos e agir em seu benefício. & 
Interface em linguagem natural (LLM) gerando 
relatórios claros para o gestor. \\ \hline
\end{tabular}

\vspace{0.1cm}
{\footnotesize Fonte – Adaptado de \citeonline{muller1999}}
\end{quadro}


De acordo com Cheng et al. \cite{cheng2026}, arquiteturas baseadas em LLMs
podem incorporar diferentes mecanismos para ampliar as capacidades dos
agentes, incluindo planejamento, memória, reflexão e utilização de
ferramentas. Esses mecanismos permitem que o agente realize tarefas que
envolvem múltiplas etapas, utilizando informações obtidas durante sua
interação com o ambiente para orientar ações posteriores. Os autores também
discutem sistemas multiagentes baseados em LLMs, nos quais diferentes agentes
podem assumir papéis distintos, comunicar-se e coordenar suas ações para a
resolução de problemas.

Dessa forma, os agentes inteligentes podem ser entendidos como sistemas cuja
característica fundamental não está apenas na execução isolada de uma tarefa,
mas na capacidade de perceber informações, processá-las, selecionar ações e
interagir com o ambiente de acordo com objetivos estabelecidos. A evolução
das arquiteturas reativas, deliberativas e híbridas para abordagens que
incorporam modelos de linguagem, memória, planejamento e utilização de
ferramentas estabelece uma ponte entre os fundamentos tradicionais dos
agentes inteligentes e as arquiteturas contemporâneas baseadas em LLMs
\cite{muller1999,cheng2026}.

\section{Large Language Models}
\label{sec:llm}

Os (\textit{Large Language Models} -- LLMs), ou Grandes Modelos de Linguagem, representam uma evolução dos modelos de processamento de linguagem natural (PLN) baseada no pré-treinamento auto-supervisionado em grandes volumes de dados textuais. Esses modelos são capazes de aprender representações linguísticas e padrões de relacionamento entre palavras, sentenças e documentos, permitindo sua aplicação em diferentes tarefas de compreensão e geração de linguagem natural. A evolução dos modelos de linguagem está diretamente relacionada ao aumento da escala dos modelos, dos conjuntos de dados utilizados no treinamento e dos recursos computacionais disponíveis \cite{zhao2026survey}.

De acordo com \citeonline{zhao2026survey}, o ciclo de vida de desenvolvimento de um LLM moderno em quatro dimensões principais: etapas distintas de treinamento: (i) o pré-treinamento, no qual o modelo adquire conhecimento geral de mundo e habilidades linguísticas por meio da predição do próximo token; (ii) o pós-treinamento, que inclui o ajuste fino supervisionado (\textit{SFT}) e o aprendizado por reforço baseado em retroalimentação humana (\textit{RLHF}); (iii) a utilização, que são as estratégias eficazes de implantação e aplicações práticas; e (iv) avaliação, que é a mensuração completa das capacidades e dos riscos. Essa perspectiva evidencia que o desempenho de um modelo não depende exclusivamente de sua arquitetura ou quantidade de parâmetros, mas também das estratégias utilizadas para treiná-lo, adaptá-lo, integrá-lo a outras fontes de informação e avaliar seus resultados \cite{zhao2026survey}.


\subsection{Retrieval-Augmented Generation}
\label{subsec:rag}

Uma das abordagens utilizadas para reduzir as limitações relacionadas ao conhecimento paramétrico dos LLMs é o \textit{Retrieval-Augmented Generation} (RAG). A abordagem combina um modelo de linguagem com mecanismos externos de recuperação de informação, permitindo que o modelo utilize documentos ou outras fontes de conhecimento durante a geração da resposta \cite{lewis2020rag}.

Em uma arquitetura RAG, o processo pode ser dividido, de forma geral, em três etapas. Na primeira etapa, denominada recuperação (\textit{retrieval}), são identificados documentos ou trechos de documentos relevantes para a consulta realizada pelo usuário. Essa recuperação pode utilizar métodos baseados em correspondência lexical ou métodos semânticos fundamentados em representações vetoriais. Na segunda etapa, as informações recuperadas são incorporadas ao contexto fornecido ao modelo de linguagem, juntamente com as instruções necessárias para a realização da tarefa. Finalmente, na terceira etapa, o LLM utiliza esse contexto enriquecido para gerar a resposta \cite{zhao2026survey}.

Entretanto, a qualidade de um sistema RAG depende diretamente da qualidade das informações recuperadas. Documentos irrelevantes, incompletos ou contraditórios podem ser incorporados ao contexto e influenciar negativamente a resposta gerada. Por esse motivo, estratégias como reranking, expansão ou reescrita de consultas, filtragem de documentos e recuperação iterativa podem ser empregadas para aumentar a relevância do contexto disponibilizado ao LLM \cite{zhao2026survey}.

\subsection{Sistemas baseados em agentes}
\label{subsec:agentic}

Uma evolução adicional na utilização dos LLMs consiste na construção de sistemas baseados em agentes (\textit{agentic systems}). Nessa abordagem, o modelo de linguagem deixa de atuar apenas como um mecanismo de entrada e saída textual e passa a desempenhar o papel de componente central de raciocínio em um sistema capaz de perceber informações, planejar ações, utilizar ferramentas e interagir com um ambiente \cite{zhao2026survey}.

Sistemas agentic podem ser estruturados a partir de três componentes funcionais principais: memória, planejamento e execução. A memória permite armazenar informações relacionadas às interações anteriores e ao contexto da tarefa. O planejamento possibilita decompor objetivos complexos em uma sequência de subtarefas e revisar estratégias de acordo com os resultados obtidos. A execução é responsável pela realização das ações planejadas, podendo envolver a utilização de ferramentas externas, APIs, mecanismos de busca, bancos de dados ou interpretadores de código \cite{zhao2026survey}.

O funcionamento desses sistemas ocorre de maneira iterativa. Inicialmente, o agente observa o estado atual da tarefa e recupera informações relevantes de sua memória. Em seguida, o componente de planejamento determina uma ação ou sequência de ações. Essas ações são executadas e os resultados obtidos são utilizados como novo contexto para o próximo ciclo de decisão. Dessa forma, o sistema estabelece um ciclo contínuo de percepção, planejamento, execução e avaliação \cite{zhao2026survey}.

Essa evolução, contudo, também introduz novos desafios. Sistemas agentic precisam lidar com crescimento do contexto, gerenciamento de memória, erros acumulados durante a execução, resultados inesperados de ferramentas e necessidade de verificar as ações realizadas. Assim, a confiabilidade desses sistemas depende não apenas das capacidades do LLM, mas também da arquitetura responsável pelo gerenciamento de memória, planejamento, execução e avaliação das ações \cite{zhao2026survey}.

De início, os LLMs foram utilizados como modelos capazes de compreender e gerar linguagem, depois, mecanismos como o RAG permitiram complementar com fontes externas e só então as arquiteturas agênticas permitiram combinar os modelos com memória, planejamento e ferramentas para a execução de tarefas mais complexas. 

\section{Arquiteturas baseadas em agentes e LLMs}


Os sistemas baseados em agentes constituem uma abordagem consolidada na Ciência da Computação para a construção de software capaz de perceber informações do ambiente, tomar decisões e executar ações orientadas a objetivos \cite{muller1999}. 

Do ponto de vista arquitetural, um agente baseado em LLM pode ser compreendido como um sistema composto por um núcleo de raciocínio associado a mecanismos externos que ampliam as capacidades do modelo. O LLM fornece a capacidade de interpretar instruções, organizar informações, realizar inferências e selecionar ações, enquanto componentes auxiliares permitem manter informações contextuais, consultar fontes externas e interagir com ferramentas ou ambientes. Essa composição permite que o agente ultrapasse algumas das limitações inerentes ao conhecimento paramétrico do modelo e execute tarefas que exigem múltiplas etapas de interação \cite{cheng2026,zhao2026survey}.

\subsection{Componentes arquiteturais de agentes baseados em LLMs}

As arquiteturas contemporâneas de agentes baseados em LLMs podem ser organizadas em torno de diferentes componentes funcionais. Embora existam variações significativas entre as propostas, três elementos fundamentais podem ser destacados: memória, planejamento e execução. Esses componentes são integrados ao modelo de linguagem em um processo iterativo no qual informações do ambiente são analisadas, ações são planejadas e executadas e os resultados obtidos são utilizados para orientar as etapas seguintes \cite{cheng2026,zhao2026survey}.

A memória é responsável por armazenar informações relevantes para a continuidade da interação e para a execução de tarefas que ultrapassam uma única chamada ao modelo. Em termos gerais, pode-se distinguir entre memória de curto prazo e memória de longo prazo. A primeira está associada às informações presentes no contexto corrente do modelo, incluindo instruções, observações e resultados intermediários. A segunda envolve mecanismos externos capazes de persistir informações entre diferentes etapas ou interações, possibilitando sua recuperação quando necessário \cite{zhao2026survey}.

O planejamento constitui outro componente fundamental. Enquanto o modelo de linguagem pode produzir uma resposta diretamente a partir de uma solicitação, um agente precisa frequentemente decompor um objetivo complexo em uma sequência de subtarefas. O planejamento permite determinar quais ações devem ser realizadas, em que ordem e quais informações ou ferramentas são necessárias para alcançar o objetivo. Em arquiteturas baseadas em LLMs, esse processo pode ocorrer de maneira iterativa, permitindo que o agente revise seu plano a partir das informações obtidas durante a execução \cite{cheng2026,zhao2026survey}.

A execução, por sua vez, corresponde à realização das ações definidas durante o planejamento. Em arquiteturas mais simples, o próprio LLM pode produzir uma ação ou instrução textual. Em arquiteturas mais sofisticadas, o agente pode utilizar ferramentas externas, como mecanismos de busca, APIs, bancos de dados, interpretadores de código e outros serviços computacionais. Essa capacidade de utilização de ferramentas amplia o espaço de ações disponível ao agente e permite que ele acesse informações ou execute operações que não poderiam ser realizadas exclusivamente a partir do conhecimento armazenado nos parâmetros do modelo \cite{zhao2026survey}.

A interação entre esses componentes pode ser representada por um ciclo contínuo. Inicialmente, o agente recebe uma tarefa e observa o estado atual do ambiente. Em seguida, recupera informações relevantes de sua memória, interpreta o contexto e estabelece um plano de ação. As ações planejadas são então executadas, podendo envolver tanto o próprio modelo quanto ferramentas externas. Os resultados dessas ações retornam ao agente na forma de novas observações, que podem ser utilizadas para atualizar o contexto e a memória e, posteriormente, revisar ou gerar novas ações. Dessa forma, a arquitetura deixa de representar um fluxo linear de entrada e saída e passa a operar como um ciclo de percepção, raciocínio, ação e feedback \cite{muller1999,cheng2026,zhao2026survey}.

\subsection{Integração com fontes externas e RAG}

A utilização de fontes externas constitui um mecanismo importante para ampliar a capacidade dos agentes baseados em LLMs. Embora modelos de linguagem possuam grande quantidade de conhecimento incorporado em seus parâmetros, esse conhecimento pode ser limitado em relação a informações específicas de domínio, conteúdos privados ou informações que sofrem alterações frequentes. A abordagem de \textit{Retrieval-Augmented Generation} (RAG) foi proposta como uma forma de combinar a capacidade generativa do modelo com informações recuperadas de fontes externas \cite{lewis2020rag}.

Em uma arquitetura RAG, documentos ou registros relevantes são inicialmente recuperados a partir de uma consulta. As informações selecionadas são então incorporadas ao contexto fornecido ao modelo, que utiliza esse conhecimento adicional durante a geração da resposta \cite{lewis2020rag}. Em agentes baseados em LLMs, esse mecanismo pode ser integrado como um componente arquitetural de recuperação de conhecimento, permitindo que o agente consulte bases documentais, bancos de dados ou outras fontes externas durante a execução de suas tarefas \cite{zhao2026survey}.

A integração entre agentes e RAG é particularmente relevante quando o sistema precisa operar sobre conhecimento especializado ou informações que não fazem parte do conhecimento paramétrico do modelo. Nesse cenário, o agente pode identificar a necessidade de informações adicionais, formular ou reformular uma consulta, recuperar os conteúdos relevantes e incorporá-los ao processo de raciocínio. O mecanismo de recuperação passa, portanto, a fazer parte do ciclo de decisão do agente, em vez de constituir apenas uma etapa independente de geração de respostas \cite{lewis2020rag,zhao2026survey}.

\subsection{Arquiteturas multiagentes}

Em arquiteturas multiagentes, a coordenação pode ocorrer de maneira centralizada ou distribuída. No primeiro caso, um agente ou componente coordenador é responsável por decompor a tarefa, distribuir responsabilidades e organizar os resultados produzidos pelos demais agentes. No segundo, os próprios agentes podem interagir entre si, compartilhando informações e tomando decisões de maneira mais descentralizada. Essas diferentes formas de organização permitem explorar especialização funcional, paralelização de tarefas e mecanismos de colaboração entre agentes \cite{cheng2026,zhao2026survey}.

A arquitetura multiagente introduz, entretanto, novos desafios relacionados à coordenação, comunicação e consistência das decisões. À medida que o número de agentes aumenta, torna-se necessário controlar a troca de informações, evitar redundância de tarefas e determinar mecanismos para resolução de conflitos. Assim, a utilização de múltiplos agentes não representa apenas a replicação de um mesmo componente, mas exige mecanismos arquiteturais específicos para orquestração e colaboração \cite{cheng2026}.

\section{Dados de Sensores}
\label{sec:dados_sensores}

No domínio dos recursos hídricos, a transição do monitoramento manual convencional para sistemas automatizados de alta frequência revolucionou a governança ambiental \cite{leigh2019framework}. Sensores multiparamétricos modernos realizam leituras semi-contínuas de variáveis como turbidez (expressa em NTU) e condutividade elétrica calibrada a 25\,°C (expressa em $\mu$S/cm), registrando dados de forma automática em intervalos temporais curtos, tipicamente de hora em hora ou de acordo com a variação do nível do leito do rio \cite{leigh2019framework}. Essas variáveis atuam frequentemente como substitutos (\textit{surrogates}) em tempo real para a estimativa de poluentes de análise química complexa e demorada, tais como sedimentos suspensos totais e concentrações de fósforo e nitrato \cite{leigh2019framework}.


\section{Detecção de Anomalias}
\label{sec:deteccao_anomalias}

No contexto de sensores físico-químicos em redes hídricas, as anomalias decorrem majoritariamente de falhas técnicas instrumentais, tais como esgotamento de baterias, calibração inadequada, falhas de transmissão de dados e, especialmente, a bioincrustação (\textit{biofouling}) das sondas submersas, as quais podem mascarar ou distorcer os padrões ecológicos reais \cite{leigh2019framework}.

No domínio dos sistemas inteligentes de monitoramento hídrico e Internet das Coisas (IoT), a necessidade de processamento contínuo em tempo real de grandes volumes de dados de sensores (\textit{data streams}) impulsionou a adoção de paradigmas de Processamento de Eventos Complexos (CEP -- \textit{Complex Event Processing}) integrados a algoritmos eficientes de aprendizado de máquina \cite{sun2019building}. Conforme demonstrado por \citeonline{sun2019building}, motores de CEP de alto desempenho (como o Apache Kafka) integrados a bancos de dados de séries temporais (como o InfluxDB) viabilizam a ingestão e o processamento de dados contínuos com baixíssima latência.

Desafios práticos da detecção de anomalias em geociências e monitoramento ambiental envolvem lidar com o severo desbalanceamento de classes e a alta sazonalidade dos ecossistemas \cite{leigh2019framework, sun2019building}. O principal desafio analítico reside na diferenciação precisa entre anomalias técnicas instrumentais e variações ecológicas extremas legítimas do próprio corpo d'água (como picos de turbidez devidos a enxurradas sazonais decorrentes de precipitação intensa), o que demanda a incorporação de conhecimento de domínio e a combinação híbrida de múltiplos métodos de detecção de forma modular e adaptável \cite{leigh2019framework, sun2019building}.

\section{Considerações Finais}
\label{sec:conclusao_fundamentacao}

A trajetória percorrida partiu dos fundamentos de Inteligência Artificial e arquiteturas de agentes inteligentes \cite{muller1999, russell2021artificial}, evoluindo para o uso de Grandes Modelos de Linguagem (\textit{Large Language Models} -- LLMs) como núcleos cognitivos de raciocínio e deliberação \cite{cheng2026, zhao2026survey}. Para mitigar as limitações de conhecimento paramétrico e alucinações dos LLMs, articulou-se a integração com a arquitetura de Geração Aumentada por Recuperação (RAG) \cite{lewis2020rag} e com ferramentas analíticas externas, como os algoritmos de detecção de anomalias em séries temporais de sensores hídricos \cite{leigh2019framework, sun2019building}. 
\chapter{Revisão Sistemática}
\chapter{Metodologia}
\chapter{Resultados experimentais}
\chapter{Conclusões}



\postextual
% ----------------------------------------------------------

% ----------------------------------------------------------
% Referências bibliográficas
% ----------------------------------------------------------
\bibliography{referencias}

\begin{apendicesenv}

\partapendices

\chapter{Protocolo da revisão de escopo}
\label{ap:protocolo}

\section{Objetivo e tipo de revisão}

Esta pesquisa adota uma revisão de escopo (\textit{scoping review}) para
mapear a natureza, a extensão e as características das evidências
disponíveis sobre o uso de modelos de linguagem de grande porte
(\textit{Large Language Models} -- LLMs), inteligência artificial e
mecanismos de apoio à decisão no contexto do monitoramento e da gestão
da qualidade da água. A revisão busca identificar abordagens, tipos de
dados, aplicações, formas de interação com usuários e lacunas de pesquisa.
O relato do processo será organizado considerando as recomendações
aplicáveis do PRISMA-ScR \cite{tricco2018prismascr}.

A revisão de escopo tem como finalidade mapear a literatura, não presumir
antecipadamente a eficácia de uma intervenção nem substituir a avaliação
individual dos estudos. Por isso, as decisões de triagem por título e
metadados são preliminares: os trabalhos mantidos deverão passar por
leitura integral antes da inclusão definitiva na síntese.

\section{Questões de pesquisa}

As questões foram definidas para mapear a contribuição da literatura em quatro
camadas complementares: aplicações de IA generativa, integração com dados de
monitoramento, apoio à decisão e limitações das soluções propostas. A análise
será orientada por categorias de extração definidas a priori e refinadas durante
a leitura integral, preservando-se categorias emergentes quando os estudos
apresentarem aspectos não previstos no protocolo.

\subsection{QP1 -- Aplicações de LLMs e IA generativa}

\textbf{Questão.} Quais aplicações de LLMs, IA generativa, geração aumentada
por recuperação (RAG) ou agentes inteligentes têm sido propostas para o
monitoramento, a interpretação ou a gestão da qualidade da água e de sistemas
hídricos?

\textbf{Motivação.} Identificar os problemas hídricos que motivam o uso dessas
tecnologias e distinguir aplicações efetivas de propostas conceituais ou usos
genéricos de IA.

\textbf{Resposta esperada.} Um mapa de finalidades e tarefas, como interpretação
de dados, consulta em linguagem natural, diagnóstico, previsão, geração de
relatórios, recomendação e coordenação de agentes, indicando o contexto hídrico
quando informado.

\textbf{Estratégia de análise.} Extrair finalidade, contexto de aplicação,
tecnologia declarada e papel desempenhado pelo LLM/agente. Agrupar os estudos
por tarefa principal e registrar separadamente trabalhos conceituais, revisões
e implementações avaliadas empiricamente.

\subsection{QP2 -- Dados de sensores e monitoramento}

\textbf{Questão.} Quais tipos de dados de sensores, IoT ou redes de
monitoramento contínuo são utilizados, e como são tratados, analisados ou
interpretados?

\textbf{Motivação.} A utilidade de sistemas inteligentes para monitoramento
hídrico depende dos dados disponíveis, da frequência de coleta e da forma como
medições e informações contextuais são integradas.

\textbf{Resposta esperada.} Uma caracterização dos dados e fontes utilizados,
como parâmetros físico-químicos, séries temporais, dados operacionais,
observações remotas ou documentos técnicos, além das etapas de processamento
relatadas.

\textbf{Estratégia de análise.} Registrar sensores/fontes, parâmetros,
frequência ou temporalidade, pré-processamento, integração de dados e métodos
analíticos. Quando um aspecto não for informado, classificá-lo como não
reportado, sem inferir sua utilização.

\subsection{QP3 -- Mecanismos de apoio à decisão}

\textbf{Questão.} Quais mecanismos de apoio à decisão são descritos, incluindo
alertas, diagnóstico de anomalias, explicações, recomendações, visualizações e
interfaces humano--IA?

\textbf{Motivação.} O objetivo aplicado da pesquisa não se limita a identificar
modelos de IA: importa compreender como os resultados são transformados em
informações utilizáveis por operadores, especialistas ou gestores.

\textbf{Resposta esperada.} Uma comparação dos mecanismos de apoio à decisão,
dos usuários a que se destinam e da participação humana prevista no fluxo de
trabalho.

\textbf{Estratégia de análise.} Codificar o tipo de saída, o usuário-alvo, o
grau de automação, a possibilidade de intervenção humana e a forma de avaliação
do apoio à decisão, quando descritos no texto.

\subsection{QP4 -- Limitações e lacunas}

\textbf{Questão.} Quais limitações técnicas, metodológicas e de avaliação são
relatadas, e quais oportunidades de pesquisa permanecem em aberto?

\textbf{Motivação.} A identificação de limitações ajuda a evitar conclusões
baseadas apenas na novidade tecnológica e orienta a definição de requisitos
para arquiteturas inteligentes aplicáveis ao monitoramento da água.

\textbf{Resposta esperada.} Uma síntese das limitações relacionadas, por
exemplo, à qualidade e disponibilidade dos dados, confiabilidade das respostas,
validação empírica, explicabilidade, integração com sistemas existentes,
segurança e supervisão humana, apenas quando sustentadas pelos estudos.

\textbf{Estratégia de análise.} Extrair limitações e recomendações declaradas
pelos autores e agrupá-las por tema. Distinguir explicitamente limitações
relatadas pelos estudos de lacunas interpretadas na síntese desta revisão.

\section{Fontes de informação e datas das buscas}

As buscas foram realizadas em cinco fontes bibliográficas: Scopus,
Web of Science (WoS), ScienceDirect, IEEE Xplore e ACM Digital Library,
em 11 de outubro de 2026. Foram recuperados 135 registros brutos, assim
distribuídos: 20 na Scopus, 23 na WoS, 39 na ScienceDirect, 1 na IEEE
Xplore e 52 na ACM Digital Library. Esses valores correspondem aos
resultados registrados para a etapa de identificação da revisão.

\section{Estratégia de busca documentada}

As consultas utilizadas em cada fonte são apresentadas a seguir,
respeitando a sintaxe e os campos de busca de cada plataforma. Não foram
aplicados filtros adicionais nas consultas descritas.

\begin{description}
    \item[IEEE Xplore]\mbox{}\\
    {\small\ttfamily\raggedright
    "water quality monitoring" AND "decision support" AND "sensor data" AND "large language models"\par}
    Sem filtros.

    \item[ScienceDirect]\mbox{}\\
    {\small\ttfamily\raggedright
    "water quality monitoring" AND "decision support" AND "sensor data" AND "large language models"\par}
    Sem filtros.

    \item[Web of Science]\mbox{}\\
    {\small\ttfamily\raggedright
    "water quality monitoring" (Topic) AND "decision support" (Topic) AND "sensor data" (Topic) AND "large language models" (Topic)\par}
    Sem filtros.

    \item[Scopus]\mbox{}\\
    {\small\ttfamily\raggedright
    "water quality monitoring" AND "decision support" AND "sensor data" AND "large language models"\par}
    Campo: \textit{All fields}.

    \item[ACM Digital Library]\mbox{}\\
    {\small\ttfamily\raggedright
    [All: "water quality monitoring"] AND [All: "decision support"] AND [All: "sensor data"] AND [All: "large language models"]\par}
    Sem filtros.
\end{description}

As expressões combinam quatro conceitos escolhidos por corresponderem ao
problema aplicado (monitoramento da qualidade da água), à finalidade decisória
(apoio à decisão), à fonte de dados de interesse (dados de sensores) e à família
de tecnologia investigada (modelos de linguagem de grande porte). O operador
AND foi usado para exigir a presença dos quatro conceitos na recuperação. Essa
opção aumenta a especificidade, mas pode reduzir a sensibilidade e deixar de
recuperar trabalhos que usem sinônimos ou não mencionem todos os conceitos nos
campos pesquisados. A sintaxe e os campossão similares mas variam por
plataforma e portanto são documentados individualmente para favorecer a transparência e
a reprodutibilidade.
\section{Critérios de elegibilidade}

Os critérios abaixo constituem a regra de triagem proposta para esta
versão e deverão ser aplicados consistentemente na leitura dos textos completos.

\subsection{Critérios de inclusão}

\begin{itemize}
    \item Publicações individuais potencialmente relacionadas a pelo menos
    uma questão de pesquisa.
    \item Artigos de periódicos e artigos individuais de conferências, incluindo
    estudos primários e revisões pertinentes, com publicação formal e revisão
    por pares verificável.
    \item Estudos sobre monitoramento ou qualidade da água; dados de
    sensores, IoT ou monitoramento contínuo; LLMs, IA generativa, RAG
    ou agentes; ou mecanismos de apoio à decisão em sistemas hídricos.
    \item Publicações com informação suficiente para uma possível
    contribuição. Primeiramente sendo feita uma consulta ao resumo e, quando elegíveis, ao texto
    completo.
\end{itemize}

\subsection{Critérios de exclusão}

\begin{itemize}
    \item Registros duplicados, preservando-se uma ocorrência por publicação.
    \item Registros de patentes, conjuntos de dados, índices de livros ou
    volumes inteiros de anais, quando não correspondam a um estudo
    individual.
    \item Livros completos e capítulos de livros, pois não serão considerados
    parte elegível desta pesquisa. A unidade de análise será o artigo
    individual publicado em periódico ou conferência com revisão por pares
    verificável.
    \item Publicações cujo conteúdo não tenha relação suficientemente
    direta com as questões de pesquisa.
    \item Manuscritos não publicados, preprints ou outros registros sem
    publicação formal verificável até a data da busca, caso não atendam ao
    recorte de literatura publicada adotado nesta revisão.
    \item Publicações para as quais, após verificação, não seja possível
    confirmar revisão por pares, conforme o recorte adotado para o corpus.
    A ausência inicial dessa informação será tratada como pendência de
    verificação, não como prova automática de ausência de revisão por pares.
    \item Texto completo indisponível somente após tentativa documentada
    de obtenção.
\end{itemize}

A ausência de metadados suficientes sobre a revisão por pares será
registrada como limitação da triagem bibliográfica. Como não haverá leitura
integral nesta etapa, os resultados serão descritos como seleção preliminar
por título e metadados, sem alegar confirmação de elegibilidade em texto completo.

\section{Procedimento de seleção}

A seleção foi realizada por identificação dos registros nas cinco fontes,
deduplicação e triagem bibliográfica por tipo documental, título e metadados.
Dos 30 registros submetidos à avaliação complementar de pertinência, 13 foram
excluídos por incompatibilidade temática. Dos 17 restantes, dois capítulos de
livro foram excluídos por não atenderem ao recorte documental adotado. Os 15
registros restantes compõem o conjunto selecionado para a análise bibliográfica
preliminar apresentada no Apêndice B. Não foi realizada leitura integral nesta
etapa, portanto, a seleção não deve ser interpretada como validação de todos
os detalhes metodológicos ou de elegibilidade a partir do texto completo.

\section{Estratégia de snowballing e avaliação da qualidade}

Foi prevista a consulta às listas de referências dos estudos mantidos após a
triagem bibliográfica, como procedimento complementar para identificar
publicações potencialmente relevantes que não tivessem sido recuperadas pelas
buscas nas bases. A execução e o resultado desse procedimento são relatados no
Apêndice B.

Não foi definido um instrumento formal de avaliação da qualidade metodológica
para esta etapa, pois o objetivo imediato é mapear a literatura e descrever suas
características. Isso não significa que a qualidade dos estudos tenha sido
considerada equivalente: limitações de método, validação e evidência empírica
serão extraídas e descritas quando a leitura integral for realizada.

\section{Extração e síntese dos dados}

Para cada estudo incluído após leitura integral, serão extraídos, quando
disponíveis: referência bibliográfica, ano, veículo de publicação,
contexto hídrico, objetivo, tipos de dados, sensores ou parâmetros
monitorados, método de processamento, papel dos LLMs/agentes, mecanismo
de apoio à decisão, forma de avaliação, resultados, limitações e lacunas.
A síntese combinará estatísticas descritivas e categorização qualitativa
orientada pelas questões de pesquisa.

\section{Limitações do registro metodológico}

Quando o ano de fascículo for posterior à data da busca, será considerada
a data de disponibilização eletrônica da publicação. Assim, a indicação
de um ano futuro no fascículo não implica, por si só, que o registro ainda
não estivesse publicado on-line na data da busca.


\chapter{Resultados da condução da revisão do estado da arte}
\label{ap:resultados}

\section{Relato da condução}

\subsection{Dados gerais da execução}

As buscas foram realizadas em cinco bases de dados científicas
no período de 4 de agosto de 2026 a 11 de outubro de 2026. A
estratégia utilizou termos relacionados ao monitoramento da qualidade
da água, apoio à decisão, dados de sensores e modelos de linguagem
de grande porte.

A string utilizada na maior parte das plataformas foi:

\begin{verbatim}
"water quality monitoring" AND "decision support"
AND "sensor data" AND "large language models"
\end{verbatim}

Não foi utilizado o termo de exclusão \textit{medical}. A data
exata de execução de cada consulta não foi preservada no registro
disponível; portanto, informa-se apenas o período geral conhecido.

A Tabela~\ref{tab:resultados-busca} apresenta a quantidade de registros
informada para cada base de dados.

\begin{table}[htbp]
    \centering
    \caption{Registros retornados pelas bases de dados}
    \label{tab:resultados-busca}
    \begin{tabular}{|l|r|}
        \hline
        \textbf{Base de dados} & \textbf{Registros} \\
        \hline
        Scopus & 15 \\
        \hline
        Web of Science & 22 \\
        \hline
        ScienceDirect & 22 \\
        \hline
        IEEE Xplore & 2 \\
        \hline
        ACM Digital Library & 1 \\
        \hline
        \textbf{Total bruto} & \textbf{62} \\
        \hline
    \end{tabular}
    % -- o\fonte nao ta definido ainda... pensar
    % \fonte{Elaboração própria, com base nos resultados das buscas.}
\end{table}

A soma dos resultados informados pelas cinco bases totalizou
62 registros brutos. Esse total representa a quantidade de resultados
retornados antes da consolidação e da remoção de duplicatas. Assim,
não corresponde necessariamente à quantidade de estudos únicos.

\subsection{Registro das consultas e limitações documentais}

As buscas foram conduzidas ao longo do período informado, mas
não há registro disponível da data exata de cada consulta. Por
esse motivo, não é possível associar uma data individual a cada
base sem introduzir uma informação não documentada.

A string principal foi utilizada na maior parte das plataformas.
As adaptações específicas de sintaxe, os campos pesquisados e
os filtros próprios de cada mecanismo não estão integralmente
documentados nesta versão. Essa limitação deverá ser considerada
na avaliação da reprodutibilidade do processo de busca.

Os endereços das plataformas foram utilizados para consultar
os resultados. Entretanto, os endereços disponíveis para IEEE
Xplore e ACM Digital Library correspondem às páginas gerais
das plataformas, não constituindo, isoladamente, evidência
persistente da consulta específica realizada.

\subsection{Processo de seleção dos estudos primários}

Os 62 registros identificados constituem o conjunto inicial de
resultados da revisão. A partir desse conjunto, devem ser realizadas
a consolidação dos registros, a identificação de duplicatas e a
avaliação da elegibilidade das publicações.

Nesta versão, as quantidades de duplicatas, exclusões por título
e resumo, textos completos avaliados e estudos incluídos ainda
não foram consolidadas. Esses valores não são apresentados como
resultados porque dependem da execução e do registro das etapas
de seleção.

A Figura~\ref{fig:fluxo-revisao} representa, de forma provisória,
a etapa inicial documentada do processo.

\begin{figure}[htbp]
    \centering
    \fbox{
        \parbox{0.85\textwidth}{
            \centering
            Identificação nas cinco bases de dados

            \vspace{0.3cm}

            Scopus: 15 registros

            Web of Science: 22 registros

            ScienceDirect: 22 registros

            IEEE Xplore: 2 registros

            ACM Digital Library: 1 registro

            \vspace{0.3cm}

            Total inicial: 62 registros

            \vspace{0.3cm}

            Próximas etapas a concluir: remoção de duplicatas,
            triagem por título e resumo, avaliação dos textos
            completos e definição do conjunto final de estudos.
        }
    }
    \caption{Etapa inicial documentada do processo de seleção}
    \label{fig:fluxo-revisao}
    % -- o\fonte nao ta definido ainda... pensar
    %\fonte{Elaboração própria.}
\end{figure}

\section{Tabela dos estudos primários selecionados}

A tabela definitiva dos estudos primários deverá ser preenchida
após a conclusão da seleção. Como essa etapa ainda não foi
consolidada, não são apresentados títulos, autores ou anos como
se constituíssem o conjunto final de estudos incluídos.

Após a seleção, a tabela deverá conter, no mínimo, o título
da publicação, os autores e o ano de publicação. Recomenda-se
também registrar a base de origem e um identificador persistente,
quando disponível.

\section{Análise descritiva}

A análise descritiva será realizada sobre o conjunto final de
estudos selecionados. Como a seleção ainda não foi concluída,
não são apresentados gráficos de distribuição temporal, veículos
de publicação, países, instituições ou autores, pois tais análises
exigem dados bibliográficos consolidados.

\subsection{Distribuição dos artigos no tempo}

Após a seleção dos estudos primários, será elaborado um gráfico
com a quantidade de publicações por ano. A análise buscará
identificar a evolução temporal das pesquisas relacionadas ao
tema da revisão.

\subsection{Tipos de veículos de publicação}

Os estudos selecionados serão classificados de acordo com
o tipo de veículo de publicação, como periódicos científicos
e conferências. A distribuição será apresentada em gráfico
e acompanhada de uma interpretação dos resultados.

\subsection{Distribuição por país, instituição e autores}

Serão analisadas as informações bibliográficas relativas aos
países, às instituições e aos autores dos estudos selecionados.
Os resultados serão apresentados por meio de gráficos ou tabelas
adequados à distribuição observada, com o objetivo de identificar
concentrações de publicações e padrões de colaboração.
\end{apendicesenv}

\begin{anexosenv}

\end{anexosenv}


\end{document}
